% Options for packages loaded elsewhere
% Options for packages loaded elsewhere
\PassOptionsToPackage{unicode}{hyperref}
\PassOptionsToPackage{hyphens}{url}
\PassOptionsToPackage{dvipsnames,svgnames,x11names}{xcolor}
%
\documentclass[
  12pt]{article}
\usepackage{xcolor}
\usepackage{amsmath,amssymb}
\setcounter{secnumdepth}{5}
\usepackage{iftex}
\ifPDFTeX
  \usepackage[T1]{fontenc}
  \usepackage[utf8]{inputenc}
  \usepackage{textcomp} % provide euro and other symbols
\else % if luatex or xetex
  \usepackage{unicode-math} % this also loads fontspec
  \defaultfontfeatures{Scale=MatchLowercase}
  \defaultfontfeatures[\rmfamily]{Ligatures=TeX,Scale=1}
\fi
\usepackage{lmodern}
\ifPDFTeX\else
  % xetex/luatex font selection
\fi
% Use upquote if available, for straight quotes in verbatim environments
\IfFileExists{upquote.sty}{\usepackage{upquote}}{}
\IfFileExists{microtype.sty}{% use microtype if available
  \usepackage[]{microtype}
  \UseMicrotypeSet[protrusion]{basicmath} % disable protrusion for tt fonts
}{}
\makeatletter
\@ifundefined{KOMAClassName}{% if non-KOMA class
  \IfFileExists{parskip.sty}{%
    \usepackage{parskip}
  }{% else
    \setlength{\parindent}{0pt}
    \setlength{\parskip}{6pt plus 2pt minus 1pt}}
}{% if KOMA class
  \KOMAoptions{parskip=half}}
\makeatother
% Make \paragraph and \subparagraph free-standing
\makeatletter
\ifx\paragraph\undefined\else
  \let\oldparagraph\paragraph
  \renewcommand{\paragraph}{
    \@ifstar
      \xxxParagraphStar
      \xxxParagraphNoStar
  }
  \newcommand{\xxxParagraphStar}[1]{\oldparagraph*{#1}\mbox{}}
  \newcommand{\xxxParagraphNoStar}[1]{\oldparagraph{#1}\mbox{}}
\fi
\ifx\subparagraph\undefined\else
  \let\oldsubparagraph\subparagraph
  \renewcommand{\subparagraph}{
    \@ifstar
      \xxxSubParagraphStar
      \xxxSubParagraphNoStar
  }
  \newcommand{\xxxSubParagraphStar}[1]{\oldsubparagraph*{#1}\mbox{}}
  \newcommand{\xxxSubParagraphNoStar}[1]{\oldsubparagraph{#1}\mbox{}}
\fi
\makeatother


\usepackage{longtable,booktabs,array}
\usepackage{calc} % for calculating minipage widths
% Correct order of tables after \paragraph or \subparagraph
\usepackage{etoolbox}
\makeatletter
\patchcmd\longtable{\par}{\if@noskipsec\mbox{}\fi\par}{}{}
\makeatother
% Allow footnotes in longtable head/foot
\IfFileExists{footnotehyper.sty}{\usepackage{footnotehyper}}{\usepackage{footnote}}
\makesavenoteenv{longtable}
\usepackage{graphicx}
\makeatletter
\newsavebox\pandoc@box
\newcommand*\pandocbounded[1]{% scales image to fit in text height/width
  \sbox\pandoc@box{#1}%
  \Gscale@div\@tempa{\textheight}{\dimexpr\ht\pandoc@box+\dp\pandoc@box\relax}%
  \Gscale@div\@tempb{\linewidth}{\wd\pandoc@box}%
  \ifdim\@tempb\p@<\@tempa\p@\let\@tempa\@tempb\fi% select the smaller of both
  \ifdim\@tempa\p@<\p@\scalebox{\@tempa}{\usebox\pandoc@box}%
  \else\usebox{\pandoc@box}%
  \fi%
}
% Set default figure placement to htbp
\def\fps@figure{htbp}
\makeatother





\setlength{\emergencystretch}{3em} % prevent overfull lines

\providecommand{\tightlist}{%
  \setlength{\itemsep}{0pt}\setlength{\parskip}{0pt}}



 
\usepackage[]{natbib}
\bibliographystyle{biom}


\addtolength{\oddsidemargin}{-.5in}%
\addtolength{\evensidemargin}{-1in}%
\addtolength{\textwidth}{1in}%
\addtolength{\textheight}{1.7in}%
\addtolength{\topmargin}{-1in}%


\usepackage{siunitx}
\usepackage{amsmath}
\usepackage{array, makecell}

\renewcommand\cellset{\renewcommand\arraystretch{0.5}%
\setlength\extrarowheight{0pt}}

\makeatletter
\def\@seccntformat#1{\@ifundefined{#1@cntformat}%
   {\csname the#1\endcsname\space}%    default
   {\csname #1@cntformat\endcsname}}%  enable individual control
\newcommand\section@cntformat{\thesection.\space} % section-level
\makeatother

\renewcommand\thesubsection{\thesection.\arabic{subsection}}
\renewcommand\thesubsubsection{\thesection.\arabic{subsection}.\arabic{subsubsection}}
\makeatletter
\@ifpackageloaded{caption}{}{\usepackage{caption}}
\AtBeginDocument{%
\ifdefined\contentsname
  \renewcommand*\contentsname{Table of contents}
\else
  \newcommand\contentsname{Table of contents}
\fi
\ifdefined\listfigurename
  \renewcommand*\listfigurename{List of Figures}
\else
  \newcommand\listfigurename{List of Figures}
\fi
\ifdefined\listtablename
  \renewcommand*\listtablename{List of Tables}
\else
  \newcommand\listtablename{List of Tables}
\fi
\ifdefined\figurename
  \renewcommand*\figurename{Figure}
\else
  \newcommand\figurename{Figure}
\fi
\ifdefined\tablename
  \renewcommand*\tablename{Table}
\else
  \newcommand\tablename{Table}
\fi
}
\@ifpackageloaded{float}{}{\usepackage{float}}
\floatstyle{ruled}
\@ifundefined{c@chapter}{\newfloat{codelisting}{h}{lop}}{\newfloat{codelisting}{h}{lop}[chapter]}
\floatname{codelisting}{Listing}
\newcommand*\listoflistings{\listof{codelisting}{List of Listings}}
\makeatother
\makeatletter
\makeatother
\makeatletter
\@ifpackageloaded{caption}{}{\usepackage{caption}}
\@ifpackageloaded{subcaption}{}{\usepackage{subcaption}}
\makeatother
\usepackage{bookmark}
\IfFileExists{xurl.sty}{\usepackage{xurl}}{} % add URL line breaks if available
\urlstyle{same}
\hypersetup{
  pdftitle={Penalized Fair Regression for Multiple Groups in Chronic Kidney Disease},
  pdfauthor={Carter H. Nakamoto; Lucia Lushi Chen; Agata Foryciarz; Sherri Rose},
  pdfkeywords={algorithmic bias, chronic kidney disease, fair
regression, penalized regression},
  colorlinks=true,
  linkcolor={blue},
  filecolor={Maroon},
  citecolor={Blue},
  urlcolor={Blue},
  pdfcreator={LaTeX via pandoc}}



\begin{document}


\def\spacingset#1{\renewcommand{\baselinestretch}%
{#1}\small\normalsize} \spacingset{1}


%%%%%%%%%%%%%%%%%%%%%%%%%%%%%%%%%%%%%%%%%%%%%%%%%%%%%%%%%%%%%%%%%%%%%%%%%%%%%%

\title{\bf Penalized Fair Regression for Multiple Groups in Chronic
Kidney Disease}
\author{
Carter H. Nakamoto, cnaks21@stanford.edu\\
Department of Health Policy, Stanford University,\\ 615 Crothers Way
Encina Commons MC6019 Stanford, CA 94305, USA\\
Lucia Lushi Chen, \\
Psyfy Inc.\\
Agata Foryciarz, \\
Department of Epidemiology and Population Health, Stanford University\\
Sherri Rose, \\
Department of Health Policy, Stanford University\\
}
\maketitle

\bigskip
\begin{center}
\textbf{Short Running Head:} Penalized Fair Regression for Multiple Groups
\end{center}

\bigskip
\begin{abstract}
Fair regression methods have the potential to mitigate societal bias
concerns in health care, but there has been little work on penalized
fair regression when multiple groups experience such bias. We propose a
general regression framework that addresses this gap with unfairness
penalties for multiple groups. Our approach is demonstrated for binary
outcomes with true positive rate disparity penalties and can be
efficiently implemented through reduction to a cost-sensitive
classification problem. We additionally introduce novel score functions
for automatically selecting penalty weights. Our penalized fair
regression methods are empirically studied in simulations, where they
achieve a fairness-accuracy frontier beyond that of existing comparison
methods. Finally, we apply these methods to a national multi-site
primary care study of chronic kidney disease to develop a fair
classifier for end-stage renal disease. There we find substantial
improvements in fairness for multiple race and ethnicity groups who
experience societal bias in the health care system without any
appreciable loss in overall fit.
\end{abstract}

\noindent%
{\it Keywords:} algorithmic bias, chronic kidney disease, fair
regression, penalized regression
\vfill

\newpage
\spacingset{1.9} % DON'T change the spacing!


\section{INTRODUCTION}\label{introduction}

Fair regression methods modify the loss function to encode fairness
measures. These loss functions typically include a single penalty or
constraint based on a measure of unfairness for one group
\citep{calders_controlling_2013, berk_convex_2017, corbett-davies_measure_2023}.
Subsequent work has developed additional technical methods for solving
single group penalized or constrained fair regression problems
\citep{dwork_decoupled_2017, chzhen_fair_2020, lohaus_too_2020, mishler_fade_2021}
and tailored them for applications in cyberbullying
\citep{gencoglu_cyberbullying_2021} and health policy
\citep{zink_fair_2020}.

Although there is a mathematical equivalence between penalized and
constrained methods such that each penalized regression has a
corresponding constrained regression with the same solution, most
methods formulate their fair regressions as constrained optimization
problems. This overlooks that penalty-based approaches have several
advantages relative to constraint-based approaches that warrant further
development. For one, the approaches are conceptually distinct. A
constraint is a user-imposed maximum level of allowable unfairness,
while a penalty is a user-imposed tradeoff between accuracy and
fairness. These conceptions are appropriate for different applications
and their varied social and regulatory contexts. The two approaches may
also require divergent computational implementations. Much of the
constrained optimization literature involves development of original
solution algorithms
\citep{kearns_preventing_2018, komiyama_nonconvex_2018}, whereas there
are many existing algorithms for minimizing penalized loss functions
under different conditions, including gradient descent, Newton's method,
Nelder-Mead, and their variations
\citep{nelder_simplex_1965, shanno_conditioning_1970, loshchilov_decoupled_2019}.
Outside of the fairness literature, penalized methods are commonly used
in statistics applications
\citep{hoerl_ridge_1970, tibshirani_regression_1996} and fair penalty
methods can benefit from advances established in those applications.

Fair regression has the potential to reduce the impact of societal bias
in health care. However, only a small number of fair regression works
have examined realistic use cases in health care
\citep{zink_fair_2020, mcguire_improving_2021, pfohl_empirical_2021, pfohl_net_2022, wang_enhancing_2025}.
Health care applications have complex ethical considerations, many of
which are unique due to the complexity of medical systems
\citep{chen_ethical_2021, chen_algorithmic_2023, feng_fair_2025}. For
example, questions about how to handle certain patient characteristics
in clinical algorithms, such as race, are nuanced such that not all fair
regression approaches may be appropriate
\citep{cusick_algorithmic_2024, zink_race_2024}.

A major limitation of most fair regression approaches is that they focus
on fairness with respect to a single group when multiple groups
typically face unfairness in real world settings
\citep{yang_fairness_2020, chen_ethical_2021, liu_practical_2022, wang_enhancing_2025}.
The joint consideration of multiple groups in fair regression is
non-trivial, with challenges related to the optimal weighting of the
groups as well as the additional computational complexity
\citep{kearns_preventing_2018, narasimhan_approximate_2020, diana_minimax_2021}.
We refer here to multiple groups rather than sensitive attributes to
avoid the particular social and sometimes legal connotations of the
latter, reflecting the breadth of possible groups for which these
methods may be appropriate \citep{deho_should_2023}.

While limited compared to methods for single groups, there is a small
literature on fair constrained regression for multiple groups
establishing basic frameworks and efficient solution methods
\citep{kearns_preventing_2018, narasimhan_learning_2018, agarwal_fair_2019, zafar_fairness_2019, narasimhan_approximate_2020, yang_fairness_2020, diana_minimax_2021, hu_sequentially_2024, wang_enhancing_2025}.
Of particular note, \citet{agarwal_reductions_2018} proposed a method
that reduces fair binary classification to a sequence of cost-sensitive
classification problems, enabling the use of efficient off-the-shelf
classifiers for fair constrained regression. This approach is flexible
across common fairness metrics and was extended in multiple later
papers, including to continuous outcomes \citep{agarwal_fair_2019} and
to non-linear but convex constraints, including demographic parity and
equalized odds, and convex or fractional-convex losses
\citep{narasimhan_learning_2018}. Much of the constrained optimization
literature involves solving a min--max saddle-point game between a
predictor and an adversary enforcing fairness constraints, often leading
to non-convex objectives, instability, substantial hyperparameter
sensitivity, and high computational cost, especially when group sizes
are imbalanced.

In contrast, we have identified even fewer papers that implement fair
penalized regression with respect to multiple groups, further
highlighting the lesser focus on this formulation of the estimation
problem. \citet{perez-suay_fair_2017} used the Hilbert Schmidt
independence criterion for fairness penalties. This method is only
compatible with one measure of fairness, demographic parity, which is
satisfied when estimated outcomes are independent of groups. A ridge
penalty method that builds on \citet{komiyama_nonconvex_2018} to enforce
demographic parity was developed in \citet{scutari_achieving_2022}.
Their approach notably includes group identifiers as predictors, which
will not be appropriate for some estimation problems. They also proposed
an extension to other fairness metrics, including equalized odds. The
focus of these first two papers on demographic parity restricts
applications in health care, where the most pressing fairness concern
may differ from and be incompatible with demographic parity
\citep{seyyed-kalantari_underdiagnosis_2021, chen_algorithmic_2023}.
\citet{pfohl_empirical_2021} and \citet{pfohl_net_2022} implement fair
penalized regression with a single penalty term that measures unfairness
with respect to a single categorical group, capturing a subset of
multiple group fair regression problems. Lastly,
\citet{liu_practical_2022} introduced a fairness measure built on
cross-covariance operators on reproducing kernel Hilbert Spaces which
were used to construct unfairness penalties. While the approach is
flexible with respect to fairness measures, its complexity presents a
potential barrier to its implementation and no code is available.

We propose a general framework for regression with unfairness penalties
for multiple groups. Logistic regression is implemented as a
demonstration estimator for this framework, and true positive rate
disparity is used as the demonstration unfairness measure, although our
framework is broad and allows for other choices. The penalized
regression problem is solved efficiently through its conversion into a
cost-sensitive classification problem and, in turn, into a weighted
classification problem. Optimal penalties for each group's unfairness
measure are automatically selected given a user-specified tradeoff
between accuracy and fairness via a random search and score function.

Our application creates a fair penalized regressor classifying whether
high-risk individuals with chronic kidney disease (CKD) will develop end
stage renal disease (ESRD) in the next year. CKD is prevalent, affecting
an estimated 14\% of the US adult population \citep{cdc_chronic_2025}.
In severe cases, CKD progresses to ESRD, requiring either dialysis or
kidney transplant, but progress can be delayed through treatment and
management. Black, Hispanic, Native American, and Asian patients all
have higher incidence and prevalence of ESRD than white patients, and
they have worse kidney function, on average, than white patients when
initiating treatment for ESRD \citep{johansen_us_2025}. Identifying
patients at risk for ESRD could help these patients engage in more
intensive treatment and slow CKD progression, and fair regression is a
potential approach to partially address disparities in progression to
ESRD. Another key algorithm for defining CKD progression, the estimated
glomerular filtration rate (eGFR) equation, was recently revised to
remove race adjustment, so CKD may be a context where it is especially
appropriate to use a fair regression method that can increase fairness
without using race adjustments \citep{cusick_algorithmic_2024}.

This work contributes to several understudied areas of the literature.
The joint consideration of fairness with respect to multiple groups
remains novel, and we additionally derive multiple ways of
mathematically synthesizing fairness information across those groups.
Further, we introduce a method for users to choose a different penalty
for each group automatically based on their preferred fairness-accuracy
tradeoff. The application of fair regression to health care with the
contextually appropriate fairness metric of true positive rate disparity
is also new.

In Section~\ref{sec-rp}, we describe the research problem, including
establishing the variables and defining fairness for this problem. In
Section~\ref{sec-methods}, we introduce the penalized regression
approach. We define the loss function, describe score functions for
selecting penalties, and characterize the random search approach for
selecting penalties. In Section~\ref{sec-sim}, we conduct a simulation
study to examine our method's performance against other available
multi-group fair regression methods with varied data distributions and
score functions. In Section~\ref{sec-application}, we apply our methods
to the chronic kidney disease study and present performance metrics for
the resulting classifiers. We conclude in Section~\ref{sec-discussion}.

\section{RESEARCH PROBLEM}\label{sec-rp}

We seek to estimate a conditional mean with \(Y\) as a function of input
variables \(X\) in a manner that is fair with respect to each group
\(g \in G\). Each group \(g\) has \(n_g\) individuals. We will also
define a reference group \(r\), a set of individuals pre-specified to be
an appropriate comparison. Let \(\mathbb{I}_{g}\) be a random variable
indicating membership in group \(g\) and \(\mathbb{I}_{r}\) be a random
variable indicating membership in reference group \(r\). Our observed
data contain \(n\) individuals, indexed by \(i\).

\subsection{Fairness Formulation}\label{fairness-formulation}

We use true positive rate equality, based on equal opportunity
\citep{hardt_equality_2016}, which holds if the true positive rate of
our estimator for \(Y\) is the same for members of group \(g\) as it is
for a reference group \(r\). A focus on true positive rate fairness (or,
equivalently, false negative rate fairness) is appropriate for the
problem of low rates of screening and diagnosis among high-risk patients
in certain groups \citep{seyyed-kalantari_underdiagnosis_2021}. This is
especially true in cases where the benefit of a diagnosis is much
greater than the potential harm, as this implies that false negatives
are much more costly than false positives and should be the primary
fairness concern. We can express different versions of this fairness
measure for estimated outcomes and for estimated probabilities. The
estimated outcome formulation is more interpretable. However, the
estimated probability formulation is differentiable, so it can be easily
minimized, and it does not depend on a particular classification
threshold.

Let \(j\) index individuals in group \(g\), \(k\) index individuals in
reference group \(r\), and \(U\) be a measure of unfairness. For
estimated outcomes, the equal opportunity unfairness measure \(U_{EO}\)
is given by: \begin{equation}\phantomsection\label{eq-fair-yhat}{
U_{EO} = \frac{\sum_{k \in r} Y_k \hat Y_k}{\sum_{k \in r} Y_k} - \frac{\sum_{j\in g}Y_j \hat Y_j}{\sum_{j\in g}Y_j},
}\end{equation} where \(\hat{Y}\) represents predicted outcomes. For
estimated probabilities, we substitute \(\hat P\) for \(\hat Y\) as has
been done in prior literature \citep{wang_enhancing_2025}. This version
of the unfairness measure, which we call \(U_{EOP}\), is given by:
\begin{equation}\phantomsection\label{eq-fair-phat}{
U_{EOP} = \frac{\sum_{k \in r} Y_k \hat P_k}{\sum_{k \in r} Y_k} - \frac{\sum_{j\in g}Y_j \hat P_j}{\sum_{j\in g}Y_j}.
}\end{equation}

\subsection{Fairness Measures for Multiple
Groups}\label{sec-measures-groups}

There is no consensus in the literature on optimal ways to evaluate
fairness for multiple groups simultaneously, as much of the past fair
regression work characterizes unfairness for each group individually
\citep{agarwal_reductions_2018, komiyama_nonconvex_2018}. Given the
unfairness metrics in Equation~\ref{eq-fair-yhat} and
Equation~\ref{eq-fair-phat} for a given group and reference, which we
will denote \(U_g\), we study three approaches to aggregating across
groups.

\begin{enumerate}
\def\labelenumi{\arabic{enumi}.}
\item
  We propose a group population-weighted average. This will be the
  primary formulation:
  \(U_{PW} = \sum_{g \in G} n_g U_g / \sum_{g \in G} n_g\).
\item
  We also propose a group-weighted average, which will give equal weight
  to smaller and larger groups:
  \(U_{GW} = \sum_{g \in G} U_g / \sum_{g \in G} 1\).
\item
  Additionally, we study a maximum across all group-level unfairness
  measures, based on \citet{ghosh_characterizing_2022}:
  \(U_{Max} = \max_{g \in G} U_g\).
\end{enumerate}

\section{METHODS}\label{sec-methods}

Our novel penalized fair regression approach generalizes to various
fairness metrics, outcome variable types, and estimator architectures.
We focus our main demonstration of the method on logistic regression
with true positive rate disparity penalties. Web Appendix A contains
loss functions for continuous outcomes with the mean residual difference
as the unfairness metric.

\subsection{Regression with Fairness Penalties}\label{sec-penloss}

The penalty term for each group \(g\) will consist of a group-specific
penalty weight \(\lambda_g\) and a measure of unfairness for that group
\(U_g\): \(\lambda_g U_g\). In the case of logistic regression, we
derive the penalized loss by combining binary cross-entropy loss with
the unfairness measure defined in Equation~\ref{eq-fair-phat} and
penalty weights \(\lambda_g\). We additionally multiply the group
unfairness penalty terms by group size \(n_g\). The resulting loss
function is \[
L = - \sum_{i=1}^n \{ Y_i\ln(\hat P_i) + (1-Y_i)\ln(1 - \hat P_i) \} + \sum_{g \in G} \lambda_g n_g \left(\frac{\sum_{k\in r}Y_k\hat P_k}{\sum_{k\in r}Y_k} - \frac{\sum_{j\in g}Y_j\hat P_j}{\sum_{j\in g}Y_j} \right). 
\] The above loss function is convex but does not have a closed form
minimizing solution. We employ an efficient optimization method based on
the reductions approach of \citet{agarwal_reductions_2018}. In this
approach, we use cost minimization loss to reconceptualize the penalized
fair regression as a cost-sensitive classification problem, which can be
implemented as a weighted non-cost-sensitive classification problem. In
this formulation, the weights are \[
W_i = |C_i^0 - C_i^1| = \left|Y_i\left \{ 2 - \sum_{g \in G}\lambda_g n_g \left( \frac{\mathbb{I}_{ri}}{\sum_{k\in r}Y_k } - \frac{\mathbb{I}_{gi}}{\sum_{j\in g}Y_j}  \right) \right\} - 1\right|,
\] and modified outcomes are \(Y_i' = \mathbb{I}(C_i^0 > C_i^1)\). It is
then possible to perform the desired fair penalized regression with a
wide array of standard classifiers (e.g., logistic regression, support
vector machines, decision trees) that are compatible with
individual-level sample weighting. More detailed derivations are in Web
Appendix B.

\subsection{Selecting Penalty Coefficients}\label{sec-pen}

Next, we identify the appropriate unfairness penalty weights
\(\lambda_g \forall g \in G\) using a random search. A number of sets of
candidate values for \(\lambda_g\) are drawn from a pre-specified
distribution, and the penalized regression estimator is fit for each
set. The penalty weight set that yields the highest score is selected.
We propose an original score function that is the weighted average of
rescaled measures of accuracy and fairness, with a user-specified
weight, \(\alpha\), which can be thought of as a fairness score weight,
to determine the relative contribution of accuracy and fairness for the
score.

The accuracy term is a transformation of accuracy rescaled such that it
is \(1\) at a reasonable lower bound for accuracy and \(0\) at a
reasonable upper bound. The basis for the accuracy term is the share of
units for which the estimator's classification matches the outcome:
\(A = \frac1n \sum_{i=1}^n \mathbb{I}(Y_i= \hat Y_{i})\). We set the
reasonable upper bound to be the unpenalized estimator's accuracy
\(A_{UP}\) and the reasonable lower bound at \(0.5\). The corresponding
rescaled accuracy term is \(A_{RS} = (A_{UP} - A)/(A_{UP} - 0.5)\).

Similarly, the fairness term is a transformation of a fairness metric
rescaled to be \(1\) at a reasonable lower bound for fairness and \(0\)
at a reasonable upper bound. The basis for the fairness term is true
positive rate disparity, measured as in Equation~\ref{eq-fair-yhat}. For
each group, we take the maximum of the unfairness term and \(0\),
considering unfairness to be the extent to which the true positive rate
of group \(g\) is lower than that of the reference group \(r\). We
synthesize across groups using either a population-weighted average,
group-weighted average, or maximum, as described in
Section~\ref{sec-measures-groups}, and rescale by dividing by the
analagous measure for the standard unpenalized estimator. Let
\(\hat Y_{UP}\) denote the unpenalized estimator's outcome estimate. The
rescaled unfairness term in the population-weighted specification is \[
U_{RS} = \frac{\sum_{g \in G}n_g \max \left( U_g, 0\right)}{\sum_{g \in G}n_g \max \left( U_{UP,g},0\right)} = 
\frac{\sum_{g \in G}n_g \max \left( \frac{\sum_{k\in r}Y_k\hat Y_k}{\sum_{k\in r}Y_k} - \frac{\sum_{j\in g}Y_j \hat Y_{j}}{\sum_{j\in g}Y_j}, 0\right)}{\sum_{g \in G}n_g \max \left(\frac{\sum_{k\in r}Y_k\hat Y_{UP,k}}{\sum_{k\in r}Y_k} - \frac{\sum_{j\in g}Y_j\hat Y_{UP,j}}{\sum_{j\in g}Y_j},0\right)}.
\] It evaluates to \(0\) if the candidate estimator achieves \(0\)
unfairness for all groups and to \(1\) if the candidate estimator has
the same unfairness as the unpenalized estimator.

The score is the negative weighted average of rescaled accuracy and
rescaled fairness with a user-specified weight \(\alpha \in [0,1]\):
\(-s = \alpha U_{RS} + (1-\alpha) A_{RS}\). For \(\alpha = 0\),
maximizing the score function is equivalent to maximizing accuracy with
no fairness concerns as in an unpenalized regression estimator. For
\(\alpha = 1\), maximizing this score function is equivalent to
minimizing true positive rate disparities without regard for accuracy.

Our penalty weight search is a random search, which yields higher
maximum scores more efficiently than a grid search
\citep{bergstra_random_2012}. All the group penalty weight candidates
are sampled from identical independent log uniform distributions. We
favor this approach over greedy sequential approaches because it allows
for optimization with respect to multiple score functions
simultaneously. The breadth of the penalty weight candidate search space
and number of candidate draws will vary for each fair penalized
regression problem, but we include some empirical demonstrations of
different choices in our simulation study.

\subsection{Computational
Implementation}\label{computational-implementation}

Penalized regression is implemented using a custom \texttt{sklearn}
estimator with the \texttt{BaseEstimator} function, primarily to encode
the construction of penalty-based sample weights for the reduction and
implement the novel score function. The logistic regression itself is
fit with the base implementation from \texttt{sklearn}. Random search
steps in the penalty weight selection process are implemented with
cross-validation facilitated by \texttt{sklearn}'s \texttt{KFold} but
without using \texttt{sklearn}'s \texttt{RandomSearchCV} to accommodate
multiple complex score functions. Score functions have original
implementations to accommodate the inclusion of multiple groups,
reference groups, and unpenalized comparison estimators. Project code is
available in a publicly accessible GitHub repository.

\section{SIMULATION STUDY}\label{sec-sim}

Our simulation was designed such that the groups have separate data
generating processes and unfair true positive rate disparities with
respect to the reference group. We had a population of 1,000,000 per
simulation with three groups, with substantial overlap between the
groups in all settings. Simulation setting 1 has all moderately sized
groups, simulation setting 2 has one small group, and all groups are
small in simulation setting 3. The data generating process is available
in Web Appendix C.

\subsection{Simulation Implementation}\label{sec-simmethods}

We aimed to explore the performance of various score functions
constructed using each of the three group synthesis mechanisms in
Section~\ref{sec-pen} with fairness score weights
\(\alpha \in \{0.1, 0.2, 0.3, 0.4, 0.5, 0.6, 0.7, 0.8, 0.9, 0.99 \}\).
Additionally, we studied different numbers of candidate penalty weight
set draws \([10, 20, \dots, 100]\) to empirically characterize changes
in accuracy and fairness as the number of candidate sets increases.
Lastly, we investigated robustness to functional form misspecification
by omitting \(X_1\) as an input variable. Comparison estimators were a
baseline unpenalized regression and the fair reductions constrained
regression method from \citet{agarwal_reductions_2018} as implemented in
the \texttt{fairlearn} python package. The exponentiated gradient
algorithm was implemented with allowed fairness constraint violations of
\(0.1\) and \(0.01\), and the grid search algorithm was used with
constraint weights of \(0.25\), \(0.5\), and \(0.75\). Otherwise,
default hyperparameters were used.

Our simulations featured 1000 replications of the full data generation
and random penalty weight search process for each simulation setting and
number of candidate penalty weight sets per random search. We used 40
draws of candidate penalty weight sets and a correctly specified
estimator with one coefficient for each input variable as default
settings. Penalty weight candidates were all sampled from
\(\text{LogUnif}(-3,1)\). Three-fold cross validation was used for the
random penalty weight search. In each random search, we simultaneously
maximized all score function variations.

\subsection{Simulation Results}\label{simulation-results}

Our novel penalized regression methods were evaluated across several key
dimensions: behavior of unfairness penalty weights by fairness score
weights, understanding penalty weights across various evaluation
measures, effect of changes in multi-group unfairness synthesis
function, and the tradeoff between fairness and accuracy. First, we
characterized the behavior of unfairness penalty weights under the
averaged score functions. The random search yielded different unfairness
penalty weights for the three groups across the user-specified fairness
score weights (Figure~\ref{fig-1}). Unfairness penalty weights
increased, on average, as the fairness score weight \(\alpha\)
increased. Unfairness penalty weights also varied across groups and
simulation settings.

The different unfairness penalty weights yielded different fairness and
accuracy measures (Figure~\ref{fig-2}). In general, for each simulation
setting, there was some range of user-specified fairness score weights
over which there was a sharp decrease in unfairness with only modest
decreases in accuracy. For simulation setting 1, our penalized method
with a fairness score weight of \(0.5\) had an accuracy decrease of 3
percentage points and a mean population-weighted unfairness decrease of
8 percentage points relative to the baseline estimator. Our penalized
method incurred steeper accuracy losses for simulation setting 3, where
a fairness score weight of 0.5 had an accuracy decrease of 7 percentage
points and a mean population-weighted unfairness decrease of 7
percentage points relative to the baseline estimator.

Next, we examine the impact of the choice of multi-group synthesis
function on accuracy and the fairness measures. All three methods for
synthesis of unfairness measures across groups in the score function ---
population-weighted mean, group-weighted mean, and maximum --- yielded
similar results with overlapping interquartile ranges, especially for
accuracy and population-weighted mean fairness (Figure~\ref{fig-3}). The
methods only diverged in maximum unfairness, where the
population-weighted mean method underperformed and the maximum method
overperformed.

We compare the fairness and accuracy results of our penalized method to
the fair reductions constrained regression method (Figure~\ref{fig-4}).
This figure also shows the fairness-accuracy frontier, the set of
results for which there is no better performance in both fairness and
accuracy \citep{liang_algorithmic_2023}. The fairness-accuracy frontier
for our penalized method was beyond that of the comparison fair
reduction methods; the comparison method did not yield estimators that
outperformed our method's estimators in both mean fairness and accuracy.
In simulation settings 1 and 2, the penalized methods were clustered
around the fairness-accuracy frontier, while the comparison method had
more variance. For simulation setting 3, which featured smaller groups,
there was more variance across simulations for both methods, with more
replications far from the fairness-accuracy frontier. Results for
simulation performance studying variation in the number of candidate
penalty weight sets and under estimator misspecification are available
in Web Figures 1-2. Accuracy and unfairness had less variation as the
number of candidate penalty weight sets increased and misspecified
estimators generally had fairness-accuracy frontiers shifted downward.

\section{CHRONIC KIDNEY DISEASE APPLICATION}\label{sec-application}

\subsection{Data}\label{data}

We applied our methods in the American Family Cohort dataset derived
from the American Board of Family Medicine PRIME Registry
\citep{phillips_prime_2017, stanford_center_for_population_health_sciences_afc_2024}.
These primary care practice electronic health record data from across
the United States contain billing information (including diagnosis and
procedure codes), drug information, and physiological patient
measurements from laboratory tests and direct clinical observation.
Using input variables measured in 2017-2018, we predicted the occurrence
of ESRD in 2019. This outcome was defined through either an ESRD
diagnosis code or serum creatinine corresponding to eGFR under 15. eGFR
was calculated with the CKD-EPI 2021 equation, which does not adjust for
race, to more accurately reflect the true physiological state of
patients \citep{inker_new_2021}.

Predictors were selected based on prior literature
\citep{segal_machine_2020, inoguchi_simplified_2022, su_machine_2022, wang_towards_2022}.
We had several physiologic predictors. The means of diastolic and
systolic blood pressure as well as the maximum and mean values for
hemoglobin A1C and billirubin were included as predictors, as was the
minimum eGFR value. Missingness was accounted for through the inclusion
of indicators of missingness as predictors. Indicators for hypertension
emergencies, type 2 diabetes, glomerulonephritis, proteinuria, and
pyelonephritis, constructed using diagnosis codes, were included as
predictors as well. Diagnosis codes used for these measures are included
in Web Table 1. Finally, binary documented sex, year of birth, race, and
ethnicity were available in the American Family Cohort dataset. We
included documented sex and age as predictors. Race and ethnicity were
used to define the groups of interest in the construction of fairness
penalties and metrics, but we intentionally did not include these
variables as predictors. We studied fairness for three race and
ethnicity groups: Black, Hispanic non-white, and Asian, with some
patients identified as both Black and Hispanic non-white as well as
Asian and Hispanic non-white. White patients were treated as the
reference group. Two groups who face health disparities in CKD had small
sample sizes (American Indian and Alaska Native as well as Native
Hawaiian and Pacific Islander) in our dataset and were not included as
groups for fairness optimization.

We restricted our sample to patients with diastolic and systolic blood
pressure measurements in 2017-2018 and to practices with at least one
ESRD patient in 2017-2019 for data quality. To target high-risk
susceptible patients, the sample was restricted to patients at least 35
years old with at least two chronic kidney disease diagnosis codes, no
documented ESRD prior to 2019, and a maximum systolic blood pressure of
at least 140 mmHg, indicating stage 2 hypertension, in 2017-2018
\citep{hsu_elevated_2005, cusick_balancing_2025}. Finally, we required
recorded values of either race or ethnicity. The resulting sample had
67397 patients, 2804 of whom developed ESRD in 2019. Detailed sample
characteristics are described in Web Table 2.

\subsection{Implementation}\label{implementation}

We used a decision threshold of 0.15 for estimated ESRD probability,
informed by evidence from prior settings that the benefits of increasing
the intensity of CKD treatment outweigh the costs and thus that false
negatives are more harmful than false positives, but we recognize that
other thresholds could also be valid
\citep{tisdale_cost-effectiveness_2022}. Penalty weight candidates were
sampled from the distribution \(\text{LogUnif}(-4,-1)\). We used
three-fold cross validation to simultaneously identify the penalty
weight sets that maximized each score function variation described in
Section~\ref{sec-simmethods}. Comparison estimators were the same as in
our simulation with the same algorithms and hyperparameters described in
Section~\ref{sec-sim}.

\subsection{Results}\label{results}

The baseline estimator had an overall accuracy of 0.95, with a
sensitivity of 0.36 and a specificity of 0.97
(Table~\ref{tbl-baselogreg}). However, its accuracy and sensitivity were
both lower for Black and Asian patients than for the overall sample,
while the opposite was true for white patients. Accuracy was lower for
Hispanic non-white patients compared to the overall sample, but not
sensitivity, although sensitivity was lower for Hispanic non-white
patients than for white patients. Hispanic non-white patients also had
the highest prevalence of ESRD among the groups we considered.

Under the 30 different score function formulations, 6 distinct sets of
penalty weights that maximized the score functions and their
corresponding penalized regression estimators were identified
(Table~\ref{tbl-esrdresmain}). In most cases, more than one combination
of synthesis method and fairness score weight yielded identical results.
Overall, the penalized estimators had 0-4 percentage point decreases in
accuracy with 1-7 percentage point decreases in unfairness. The
penalized regression estimator with group-weighting and fairness score
weight 0.1 had the same performance as the population-weighted estimator
with fairness score weight 0.1 or 0.2. These three estimators had less
than a percentage point difference in accuracy, 1 percentage point
higher sensitivity, and 7 percentage points lower mean unfairness (a
78\% reduction) compared to the baseline, including reducing unfairness
for Black and Hispanic non-white patients to 0. Thus, these estimators
provided no appreciable loss in overall fit with only improvements in
fairness. However, across most estimators, substantial unfairness
remained for Asian patients, the smallest group we considered in our
data sample, comprising only 2\% of patients. The fair reductions
constrained regression comparison estimators either had extremely low
sensitivity or extremely low accuracy even when they improved
unfairness, which is inappropriate for this application
\citep{tisdale_cost-effectiveness_2022}.

\section{DISCUSSION}\label{sec-discussion}

We proposed a new fair penalized regression method that jointly
considers multiple groups with a flexible and efficient implementation.
Our method robustly achieved a fairness-accuracy frontier beyond that of
comparison fair regression methods across numerous simulation settings.
We applied the method to primary care data to develop a classifier for
ESRD that had a preferable balance of accuracy, sensitivity, and true
positive rate fairness across racial and ethnic groups compared to
existing methods. While the overall best fair penalized regression
method will depend on underlying (typically unknown) data structure and
application, our framework is well suited to a priori specified
benchmarks for selecting among the synthesis mechanisms and fairness
weights. It can seamlessly be incorporated into ensembles, for example.

Although this framework is general, we do not empirically demonstrate
how it performs for other fairness definitions, estimator types, or
outcomes. Future areas of work include adapting our framework for such
cases. We also found that for small groups of interest, our proposed
method may have high variance and could be outperformed by alternative
fair regression methods. In our application, we highlight that there is
no universal standard for how race is measured across clinics in the
American Family Cohort. We also lacked the sample size to study fairness
for some racial groups, including American Indian and Alaska Native and
Native Hawaiian and Pacific Islander patients. Furthermore, fair
regression is just one aspect of ethical machine learning and should be
considered alongside the full machine learning pipeline in health care
\citep{chen_ethical_2021}.

We made several novel contributions to the literature beyond the
inherent novelty of a penalized regression that is fair with respect to
multiple groups. In adapting the reductions method for penalized
regression, we detailed efficient implementation practices for fair
penalized regressions. The automatic penalty selection process we
developed includes a novel score function that rescales and synthesizes
accuracy and fairness metrics. The use of a random search for penalty
weights is novel in this literature, and we characterized our
implementation details and the random search performance under different
hyperparameters.

\subsection*{Acknowledgements}\label{acknowledgements}
\addcontentsline{toc}{subsection}{Acknowledgements}

Data for this project were accessed using the Stanford Center for
Population Health Sciences Data Core. The PHS Data Core is supported by
a National Institutes of Health National Center for Advancing
Translational Science Clinical and Translational Science Award
(UL1TR003142) and from Internal Stanford funding. The content is solely
the responsibility of the authors and does not necessarily represent the
official views of the NIH. We wish to acknowledge and thank the ABFM
PRIME Registry participating clinicians and the American Board of Family
Medicine, without whom the American Family Cohort would not be possible.
We also thank members of the Health Policy and Data Science Lab for
their feedback on this research with special thanks to Malcolm Barrett
for expert code review.

This project was supported by NIH grant R01LM013989, NSF grant
DGE-2146755, and AHRQ grant 5T32HS026128-07. The authors have no
conflicts of interest to disclose. The project falls under a Stanford
IRB, which approved human subjects involvement and waived consent.

\subsection*{Data availability
statement}\label{data-availability-statement}
\addcontentsline{toc}{subsection}{Data availability statement}

The data underlying this article cannot be shared publicly due to data
use agreement restrictions.

\subsection*{References}\label{references}
\addcontentsline{toc}{subsection}{References}

\renewcommand{\bibsection}{}
\bibliography{fairpenalties.bib}

\subsection*{Supplementary Material}\label{supplementary-material}
\addcontentsline{toc}{subsection}{Supplementary Material}

Web Appendices, Tables, and Figures, and code referenced in Sections 3,
3.1, 4, 4.2, and 5.1 are available with this paper. Code is available in
a public GitHub repository at the link
https://github.com/StanfordHPDS/FairPenalizedRegression\_MultipleGroups.

\newpage{}

\begin{figure}

\begin{minipage}{0.30\linewidth}

\centering{

\pandocbounded{\includegraphics[keepaspectratio]{penalty_writeup_files/figure-pdf/fig-1-output-1.pdf}}

}

\subcaption{\label{fig-1-1}Simulation Setting 1}

\end{minipage}%
%
\begin{minipage}{0.05\linewidth}
~\end{minipage}%
%
\begin{minipage}{0.30\linewidth}

\centering{

\pandocbounded{\includegraphics[keepaspectratio]{penalty_writeup_files/figure-pdf/fig-1-output-2.pdf}}

}

\subcaption{\label{fig-1-2}Simulation Setting 2}

\end{minipage}%
%
\begin{minipage}{0.05\linewidth}
~\end{minipage}%
%
\begin{minipage}{0.30\linewidth}

\centering{

\pandocbounded{\includegraphics[keepaspectratio]{penalty_writeup_files/figure-pdf/fig-1-output-3.pdf}}

}

\subcaption{\label{fig-1-3}Simulation Setting 3}

\end{minipage}%

\caption{\label{fig-1}\textbf{Penalty Weights for Each Group by Fairness Score Weight.}
Confidence intervals represent the empirical 25th and 75th percentiles
of replications.}

\end{figure}%

\begin{figure}

\begin{minipage}{0.47\linewidth}

\centering{

\pandocbounded{\includegraphics[keepaspectratio]{penalty_writeup_files/figure-pdf/fig-2-output-1.pdf}}

}

\subcaption{\label{fig-2-1}Accuracy}

\end{minipage}%
%
\begin{minipage}{0.06\linewidth}
~\end{minipage}%
%
\begin{minipage}{0.47\linewidth}

\centering{

\pandocbounded{\includegraphics[keepaspectratio]{penalty_writeup_files/figure-pdf/fig-2-output-2.pdf}}

}

\subcaption{\label{fig-2-2}Population-Weighted Mean Unfairness}

\end{minipage}%
\newline
\begin{minipage}{0.47\linewidth}

\centering{

\pandocbounded{\includegraphics[keepaspectratio]{penalty_writeup_files/figure-pdf/fig-2-output-3.pdf}}

}

\subcaption{\label{fig-2-3}Group-Weighted Mean Unfairness}

\end{minipage}%
%
\begin{minipage}{0.06\linewidth}
~\end{minipage}%
%
\begin{minipage}{0.47\linewidth}

\centering{

\pandocbounded{\includegraphics[keepaspectratio]{penalty_writeup_files/figure-pdf/fig-2-output-4.pdf}}

}

\subcaption{\label{fig-2-4}Max Group Unfairness}

\end{minipage}%

\caption{\label{fig-2}\textbf{Fair Regression Estimator Accuracy and Unfairness for Different Fairness Score Weights by Simulation Setting.}
Confidence intervals represent the empirical 25th and 75th percentiles
of replications.}

\end{figure}%

\begin{figure}

\begin{minipage}{0.47\linewidth}

\centering{

\pandocbounded{\includegraphics[keepaspectratio]{penalty_writeup_files/figure-pdf/fig-3-output-1.pdf}}

}

\subcaption{\label{fig-3-1}Accuracy}

\end{minipage}%
%
\begin{minipage}{0.06\linewidth}
~\end{minipage}%
%
\begin{minipage}{0.47\linewidth}

\centering{

\pandocbounded{\includegraphics[keepaspectratio]{penalty_writeup_files/figure-pdf/fig-3-output-2.pdf}}

}

\subcaption{\label{fig-3-2}Population-Weighted Mean Unfairness}

\end{minipage}%
\newline
\begin{minipage}{0.47\linewidth}

\centering{

\pandocbounded{\includegraphics[keepaspectratio]{penalty_writeup_files/figure-pdf/fig-3-output-3.pdf}}

}

\subcaption{\label{fig-3-3}Group-Weighted Mean Unfairness}

\end{minipage}%
%
\begin{minipage}{0.06\linewidth}
~\end{minipage}%
%
\begin{minipage}{0.47\linewidth}

\centering{

\pandocbounded{\includegraphics[keepaspectratio]{penalty_writeup_files/figure-pdf/fig-3-output-4.pdf}}

}

\subcaption{\label{fig-3-4}Max Group Unfairness}

\end{minipage}%

\caption{\label{fig-3}\textbf{Fair Regression Estimator Accuracy and Unfairness for Different Fairness Score Weights by Multi-Group Unfairness Synthesis Function.}
The subfigures correspond to performance measures for fitted estimators,
while colors correspond to how the fairness components of the score
functions were calculated. Confidence intervals represent the empirical
25th and 75th percentiles of replications.}

\end{figure}%

\begin{figure}

\begin{minipage}{0.30\linewidth}

\centering{

\pandocbounded{\includegraphics[keepaspectratio]{penalty_writeup_files/figure-pdf/fig-4-output-1.pdf}}

}

\subcaption{\label{fig-4-1}Simulation Setting 1}

\end{minipage}%
%
\begin{minipage}{0.05\linewidth}
~\end{minipage}%
%
\begin{minipage}{0.30\linewidth}

\centering{

\pandocbounded{\includegraphics[keepaspectratio]{penalty_writeup_files/figure-pdf/fig-4-output-2.pdf}}

}

\subcaption{\label{fig-4-2}Simulation Setting 2}

\end{minipage}%
%
\begin{minipage}{0.05\linewidth}
~\end{minipage}%
%
\begin{minipage}{0.30\linewidth}

\centering{

\pandocbounded{\includegraphics[keepaspectratio]{penalty_writeup_files/figure-pdf/fig-4-output-3.pdf}}

}

\subcaption{\label{fig-4-3}Simulation Setting 3}

\end{minipage}%

\caption{\label{fig-4}\textbf{Accuracy and Mean Fairness for Proposed and Comparison Fair Reductions Methods by Simulation Setting.}
Each plotted point corresponds to the result of a single simulation
experiment under one of the score function variations or enumerated fair
reductions hyperparameters.}

\end{figure}%

\newpage{}

\begin{table}

\caption{\label{tbl-baselogreg}\textbf{Unpenalized Regression Performance by Race and Ethnicity Group.}}

\centering{

\begin{tabular}{@{} l *{5}{S[table-format=3.2]} @{}}
\toprule
 & {n} & {Prevalence} & {Sensitivity} & {Specificity} & {Accuracy} \\
\midrule
Black & 8094 & 0.07 & 0.34 & 0.96 & 0.92 \\
Hispanic non-white & 2678 & 0.08 & 0.36 & 0.97 & 0.93 \\
Asian & 1427 & 0.07 & 0.19 & 0.98 & 0.93 \\
White & 53402 & 0.04 & 0.39 & 0.98 & 0.96 \\
Total & 67397 & 0.04 & 0.36 & 0.97 & 0.95 \\
\bottomrule
\end{tabular}

}

\end{table}%

\begin{table}

\caption{\label{tbl-esrdresmain}\textbf{Performance for Baseline Unpenalized Regression, Novel Penalized Regressions, and Fair Reductions Methods.}}

\centering{

\sisetup{detect-weight=true,detect-inline-weight=math, text-series-to-math = true, propagate-math-font = true}
\newrobustcmd\B{\DeclareFontSeriesDefault[rm]{bf}{b}\bfseries}
\addtolength{\jot}{-1.5em}
\begin{tabular}{@{} l *{6}{S[table-format=3.2]} @{}}
\toprule
\multicolumn{3}{r}{} & \multicolumn{4}{c}{Unfairness} \\
\cline{4-7}{Method} & {Accuracy} & {Sensitivity} & {Mean} & {Black} & {Asian} & {HNW} \\
\midrule
Baseline & 0.95 & 0.36 & 0.09 & 0.06 & 0.24 & 0.07 \\
\makecell[l]{Group $\alpha=0.1$\\Pop $\alpha=0.1, 0.2$} & 0.95 & 0.37 & 0.02 & 0.00 & 0.19 & 0.00 \\
\makecell[l]{Max $\alpha=0.1$} & 0.95 & 0.29 & 0.08 & 0.06 & 0.17 & 0.07 \\
\makecell[l]{Group $\alpha=0.2, 0.3$} & 0.93 & 0.38 & 0.03 & 0.00 & 0.23 & 0.00 \\
\makecell[l]{Max $\alpha=0.2, 0.3, 0.4, 0.5, 0.6$} & 0.94 & 0.31 & 0.08 & 0.07 & 0.11 & 0.08 \\
\makecell[l]{$\begin{aligned} \text{Group } \alpha=\,&0.4, 0.5, 0.6, 0.7 \\ &0.8, 0.9, 0.99 \end{aligned}$\\[-10pt]$\begin{aligned} \text{Pop } \alpha=\,&0.3, 0.4, 0.5, 0.6 \\ &0.7, 0.8, 0.9, 0.99 \end{aligned}$} & 0.92 & 0.39 & 0.02 & 0.00 & 0.17 & 0.00 \\
\makecell[l]{Max $\alpha=0.7, 0.8, 0.9, 0.99$} & 0.91 & 0.33 & 0.08 & 0.08 & 0.00 & 0.13 \\
ExpGrad, Ep = 0.1 & 0.96 & 0.15 & 0.04 & 0.03 & 0.12 & 0.02 \\
ExpGrad, Ep = 0.01 & 0.93 & 0.13 & 0.02 & 0.01 & 0.10 & 0.00 \\
Grid, CW = 0.25 & 0.32 & 0.90 & 0.01 & 0.00 & 0.12 & 0.00 \\
Grid, CW = 0.5 & 0.11 & 0.95 & 0.01 & 0.02 & 0.00 & 0.00 \\
Grid, CW = 0.75 & 0.11 & 0.95 & 0.01 & 0.02 & 0.00 & 0.00 \\
\bottomrule
\end{tabular}
\vspace{0.2cm} 

 \begin{flushleft} \small{\textbf{Note:} Novel fair penalized regressions were defined based on the multi-group unfairness synthesis mechanism: population-weighted average (Pop), group-weighted average (Group), and maximum (Max) as well as fairness score weight ($\alpha$). Fair reductions comparisons were exponentiated gradient (ExpGrad) or grid search (Grid) methods, varying constraint flexibility (Ep) or constraint weight (CW), respectively. `HNW' is Hispanic non-white.} \end{flushleft}

}

\end{table}%





\end{document}
